\documentclass[12pt, letterpaper]{article}
\usepackage{blindtext}
\usepackage[T1]{fontenc}
\usepackage[utf8]{inputenc}
\usepackage{amsfonts}
\usepackage{mathtools}
\usepackage{amsmath}
\usepackage{amssymb}
\usepackage{enumerate}
\usepackage{mathabx}
\usepackage{apacite}
\usepackage{subfiles}
\usepackage{amstext} % for \text macro
\usepackage{array}
\usepackage{amsthm}% for \newcolumntype macro
\usepackage[T1]{fontenc}
\usepackage{tabularx}
\usepackage{etoolbox}
\usepackage{placeins}
\usepackage[parfill]{parskip}
\usepackage[]{algorithm2e}
\usepackage{empheq}
\DeclarePairedDelimiter\abs{\lvert}{\rvert}
\title{Wybrane metody Kryłowa}
\author{Ryszard Eggink} 
\date{Czerwiec 2020}
\newtheoremstyle{custom}%    <name>
                {\topsep}%   <space above>
                {\topsep}%   <space below>
                {\itshape}%  <body font>
                {}%          <indent amount>
                {\bfseries}% <Theorem head font>
                {.}%         <punctuation after theorem head>
                {\newline}%  <space after theorem head> (default .5em)
                {}%          <Theorem head spec>
\newtheorem{tw}{Twierdzenie}[section]
\newtheorem{lem}{Lemat}[section]
\newtheorem*{przy*}{Przykład}
\newtheorem{stw}{Stwierdzenie}[section]
\newtheorem{wniosek}{Wniosek}
\newtheorem{alg}{Algorytm}
\newenvironment{manualtheorem}[1]{%
  \renewcommand\thealg{#1}%
  \alg
}{\endalg}
\theoremstyle{plain}
\newtheorem*{dow*}{Dowód}
\newcommand{\norm}[1]{\left\lVert#1\right\rVert}
\DeclareMathOperator\Ric{Ric} 
\DeclareMathOperator*{\argmin}{arg\,min}
\renewcommand*{\proofname}{Dowód}
\renewcommand{\abstractname}{Abstrakt}
\numberwithin{equation}{section}
\begin{document}
\maketitle
\begin{abstract}
W tej pracy analizowane są metody Kryłowa z książki "Iterative Methods for Sparse Linear Systems" \cite{saad2003iterative} i testowane na wybranym zadaniu eliptycznym.
\end{abstract}
\section{Problem testowy}
\subsection{Zadanie jednowymiarowe}
Będziemy testować algorytmy na problemach równań różniczkowych cząstkowych, najpierw jednowymiarowym, potem dwuwymiarowym. Za pomocą różnic skończonych można zamienić takie zadanie do układu równań liniowych.
Rozpatrzmy następujące zadanie:
\begin{equation}
\begin{split}
u''(x)+bu'(x)+cu(x)&=f (x)\:\:\:\forall x\in\Omega=[d,e] \\
u(d)&=\alpha\\
u(e)&=\beta
\end{split}
\end{equation}
dla ustalonej stałej $b$ oraz $c>0$, ustalonego odcinka $[d,e]$ oraz znanych wartości $\alpha$, $\beta$.
Będziemy przybliżać pierwszą i drugą pochodną za pomocą różnic skończonych:
\begin{equation}
\begin{split}
&\partial_{h}u(x)=\frac{u(x+h)-u(x)}{h}\\
&\partial\overline{\partial}_{h}u(x)=\frac{u(x-h)-2u(x)+u(x+h)}{h^2}
\end{split}
\end{equation}
Wprowadzamy siatkę dla $h=\frac{e-d}{N}$:
$$\overline{\Omega}=\{d+kh\}_{k=0}^N,$$
z $\Omega_h=\{d+kh\}_{k=1}^{N-1}$ oraz $\partial\Omega_h=\{d+kh\}_{k=\{0,N}=\{d,e\}$, możemy wówczas zdefiniować następujące zadanie: znaleźć funkcję $u_h$ określoną na siatce $\overline{\Omega}_h$ taką, że 
\begin{empheq}[left ={\empheqlbrace}]{alignat* = 2}
-\partial\overline{\partial}_{h}u_{h}(x)+b\partial_{h}u_{h}(x)+cu_{h}(x)=f(x)\:\:\:&x\in\Omega_{h}, \\
u(x)=g(x)\:\:\: &x\in\partial\Omega_{h},
\end{empheq}
gdzie $g:\partial\Omega\rightarrow\mathbb{R}$ przyjmuje wartości $g(d)=\alpha$, $g(e)=\beta$. Sprowadzimy teraz powyższe równania do postaci macierzowej. Przyjmiemy następujące oznaczenia $u_k=u_h(x_k)=u_h(d+kh)$ oraz $f(x_k)=f_k$ dostajemy następujący układ równań liniowych:
\begin{empheq}[left ={\empheqlbrace}]{alignat* = 2}
\frac{-u_{k-1}+2u_k-u_{k+1}}{h^2}+\frac{b(u_{k+1}-u_k)}{h}+cu_k&=f_k \\
u_0&=g(d)\\
u_N&=g(e)
\end{empheq}
Przekształcając do postaci macierzowej otrzymujemy:
\begin{align*}
\begin{pmatrix}
2-bh+ch^2 & bh-1& 0 & \cdots & a_{1,N-1} \\
-1 & 2-bh+ch^2&bh-1 & \cdots & a_{2,N} \\
\vdots  &\ddots& \ddots  & \ddots & \vdots  \\
0 & 0 &-1&  & 2-bh+ch^2&bh-1 
\end{pmatrix}
\begin{bmatrix}
  u_{1,1}  \\
  u_{1,2}  \\
  \vdots \\
  u_{1,N-1}  \\
  u_{2,1}  \\
  u_{2,2}  \\
  \vdots \\
  u_{N-1,N-2}  \\
  u_{N-1,N-1}  \\
\end{bmatrix} =
\begin{bmatrix}
f_{1,1}+\frac{g(d,f)}{h_1^2}+\frac{g(d,f)}{h_2^2}  \\
  f_{1,2}+\frac{g(d,f+2h_2)}{h_2^2}  \\
  \vdots \\
  f_{1,N-1}+\frac{g(d,f+(N-1)h_2)}{h_2^2}  \\
  f_{2,1}+\frac{g(d+2h_1,f)}{h_1^2}  \\
  f_{2,2}  \\
  \vdots \\
  f_{N-1,N-2}+\frac{g(e,f+(N-2)h_2)}{h_1^2}  \\
  f_{N-1,N-1}+\frac{g(e,g)}{h_1^2}+\frac{g(e,g)}{h_2^2}  \\
\end{bmatrix}
\end{align*}
\subsection{Zadanie dwuwymiarowe}
Rozważmy zadanie eliptyczne dwuwymiarowe na prostokącie $\overline{\Omega}=[d,e]\times [f,g]$. Zadanie polega na znalezieniu $u\in C^2(\Omega)\cap C(\overline{\Omega})$ takiego, że
\begin{equation}
\begin{split}
\Delta u(x)+&b^T\nabla u(x)+cu(x)=f (x)\:\:\:\forall x\in\Omega \\
&u(s)=g(s)\:\:\: s\in \partial \Omega,
\end{split}
\end{equation}
gdzie $\Delta u(x)=\sum_{k=1}^{2}\frac{\partial^2u}{\partial x^2_k}$, $b$ to dany wektor, $c>0$ jest ustaloną stałą, $f$ to funkcja ciągła na $\Omega$, a $g$ to funkcja ciągła na $\partial\Omega$.

Dla ustalonego kroku $h>0$ będziemy aproksymować pierwszą i drugą pochodną cząstkową:
\begin{equation}
\begin{split}
&\partial_{k,h}u(x)=\frac{u(x+he_k)-u(x)}{h}\\
&\partial\overline{\partial}_{k,h}u(x)=\frac{u(x-he_k)-2u(x)+u(x+he_k)}{h^2},\:\:\: k=1,2,
\end{split}
\end{equation}
gdzie $e_k$ to $k-ty$ wektor jednostkowy.
Wprowadzamy siatkę dla $h_1=\frac{e-d}{N}$, $h_2=\frac{g-f}{N}$:
\begin{equation}
\overline{\Omega}_{h_1,h_2}=\{(d+kh_1,f+lh_2)\}_{k,l=0}^{N}
\end{equation}
i jej odpowiednie podzbiory $\Omega_{h_1,h_2}=\{(d+kh_1,f+lh_2)\}_{k,l=1}^{N-1}\subset \Omega$, $\partial\Omega_{h_1,h_2}=\{(d+kh_1,f+lh_2)\}_{k=\{0,N\},\: l=0,\dots, N, l=\{0,N\},\: k=0,\dots,N}\subset\partial\Omega$.
Definiujemy następujące zadanie dyskretne: chcemy znaleźć funkcję $u_{h_1,h_2}$ określoną na siatce $\overline{\Omega}_{h_1,h_2}$ taką, że
\begin{empheq}[left ={\empheqlbrace}]{alignat* = 2}
-\sum_{k=1,2}\partial\overline{\partial}_{k,h_k}u_{h_1,h_2}(x)+\sum_{k=1,2}b_k\partial_{k,h_k}u_{h_1,h_2}(x)+cu_{h_1,h_2}(x)=f(x)\:\:\:&x\in\Omega_{h_1,h_2}, \\
u(x)=g(x)\:\:\: &x\in\partial\Omega_{h_1,h_2}
\end{empheq}
Oznaczmy dla uproszczenia zapisu $u_{k,l}=u_{h_1,h_2}(x_{k,l})=u_{h_1,h_2}(d+kh_1,f+lh_2)$ oraz $f_{k,l}=f(x_{k,l})=f(d+kh_1,f+lh_2)$. Otrzymujemy wówczas następujący układ liniowy.
\begin{empheq}[left ={\empheqlbrace}]{alignat* = 6}
&\frac{-u_{k+1,l}+2u_{k,l}-u_{k-1,l}}{h^2_1}+\frac{-u_{k,l+1}+2u_{k,l}-u_{k,l-1}}{h^2_2}+\frac{u_{k+1,l}-u_{k,l}}{h_1}+\frac{u_{k,l+1}-u_{k,l}}{h_2}+\\
&+cu_{k,l}=f_{k,l},\:\:\: k,l=1,\dots,N-1\\
\:\:&u_{0,l}=g(d,f+lh_2),\:\:\: l=0,\dots,N\\
\:\:&u_{k,N}=g(d+kh_1,g),\:\:\: k=0,\dots,N\\
\:\:&u_{N,l}=g(e,f+lh_2),\:\:\: l=0,\dots,N\\
\:\:&u_{k,0}=g(d+kh_1,f),\:\:\: k=0,\dots,N
\end{empheq}
Co można zapisać w postaci macierzowej podstawiając warunki brzegowe jako
\begin{align*}
\begin{pmatrix}
a_{1,1} & a_{1,2} & \cdots & a_{1,N-1} \\
a_{2,1} & a_{2,2} & \cdots & a_{2,N} \\
\vdots  & \vdots  & \ddots & \vdots  \\
a_{N-1,1} & a_{N-1,2} & \cdots & a_{N-1,N-1} 
\end{pmatrix}
\begin{bmatrix}
  u_{1,1}  \\
  u_{1,2}  \\
  \vdots \\
  u_{1,N-1}  \\
  u_{2,1}  \\
  u_{2,2}  \\
  \vdots \\
  u_{N-1,N-2}  \\
  u_{N-1,N-1}  \\
\end{bmatrix} =
\begin{bmatrix}
f_{1,1}+\frac{g(d,f)}{h_1^2}+\frac{g(d,f)}{h_2^2}  \\
  f_{1,2}+\frac{g(d,f+2h_2)}{h_2^2}  \\
  \vdots \\
  f_{1,N-1}+\frac{g(d,f+(N-1)h_2)}{h_2^2}  \\
  f_{2,1}+\frac{g(d+2h_1,f)}{h_1^2}  \\
  f_{2,2}  \\
  \vdots \\
  f_{N-1,N-2}+\frac{g(e,f+(N-2)h_2)}{h_1^2}  \\
  f_{N-1,N-1}+\frac{g(e,g)}{h_1^2}+\frac{g(e,g)}{h_2^2}  \\
\end{bmatrix}
\end{align*}
\bibliographystyle{apacite}
\bibliography{reference}
\end{document}